% =============================================================================
% Chapter 12: E-R Μοντέλο - Λυμένες Ασκήσεις
% Data Mining Study Guide - NTUA
% =============================================================================

\chapter{E-R Μοντέλο και Σχεδιασμός Βάσεων Δεδομένων - Λυμένες Ασκήσεις}
\label{ch:er-model-exercises}

Αυτό το κεφάλαιο περιέχει λυμένες ασκήσεις για το E-R μοντέλο και τον σχεδιασμό
βάσεων δεδομένων.

% -----------------------------------------------------------------------------
% Exercise 1.1
% -----------------------------------------------------------------------------
\section{Άσκηση 1.1: Σχεδιασμός E-R από Απαιτήσεις}

\begin{formulabox}[Εκφώνηση]
Σχεδιάστε ένα E-R διάγραμμα για ένα σύστημα music streaming με τις εξής απαιτήσεις:
\begin{itemize}
    \item Κάθε \textbf{τραγούδι} έχει μοναδικό ID, τίτλο, διάρκεια και έτος κυκλοφορίας
    \item Κάθε \textbf{καλλιτέχνης} έχει μοναδικό ID, όνομα και χώρα προέλευσης
    \item Κάθε τραγούδι ανήκει σε ακριβώς έναν καλλιτέχνη, αλλά ένας καλλιτέχνης μπορεί να έχει πολλά τραγούδια
    \item Οι \textbf{χρήστες} έχουν ID, username και email
    \item Οι χρήστες δημιουργούν \textbf{playlists} με όνομα και ημερομηνία δημιουργίας
    \item Μια playlist μπορεί να περιέχει πολλά τραγούδια και ένα τραγούδι μπορεί να είναι σε πολλές playlists
\end{itemize}
\end{formulabox}

\subsection*{Λύση}

\textbf{Βήμα 1: Αναγνώριση Οντοτήτων}

\textit{Γιατί:} Πρώτα εντοπίζουμε τα «αντικείμενα» του συστήματος που έχουν ανεξάρτητη ύπαρξη.

Οντότητες:
\begin{itemize}
    \item \texttt{Song} (τραγούδι)
    \item \texttt{Artist} (καλλιτέχνης)
    \item \texttt{User} (χρήστης)
    \item \texttt{Playlist}
\end{itemize}

\textbf{Βήμα 2: Αναγνώριση Γνωρισμάτων}

\begin{center}
\begin{tabular}{ll}
\toprule
\textbf{Οντότητα} & \textbf{Γνωρίσματα} \\
\midrule
Song & \underline{song\_id}, title, duration, release\_year \\
Artist & \underline{artist\_id}, name, country \\
User & \underline{user\_id}, username, email \\
Playlist & \underline{playlist\_id}, name, created\_date \\
\bottomrule
\end{tabular}
\end{center}

\textbf{Βήμα 3: Αναγνώριση Συσχετίσεων και Cardinality}

\begin{itemize}
    \item \texttt{Artist} --- \texttt{creates} --- \texttt{Song}: \textbf{1:N}
    \item \texttt{User} --- \texttt{owns} --- \texttt{Playlist}: \textbf{1:N}
    \item \texttt{Playlist} --- \texttt{contains} --- \texttt{Song}: \textbf{M:N}
\end{itemize}

\textbf{Βήμα 4: E-R Διάγραμμα}

\begin{center}
\begin{tikzpicture}[node distance=3cm,
    entity/.style={rectangle, draw=primary, fill=box-blue, minimum width=2cm, minimum height=1cm, font=\bfseries},
    relationship/.style={diamond, draw=secondary, fill=box-yellow, aspect=2, font=\small},
    attribute/.style={ellipse, draw=gray, fill=light-gray, font=\small\ttfamily, inner sep=2pt}]

    % Entities
    \node[entity] (artist) {Artist};
    \node[entity, right=4cm of artist] (song) {Song};
    \node[entity, below=3cm of artist] (user) {User};
    \node[entity, right=4cm of user] (playlist) {Playlist};

    % Relationships
    \node[relationship] (creates) at ($(artist)!0.5!(song)$) {creates};
    \node[relationship] (owns) at ($(user)!0.5!(playlist)$) {owns};
    \node[relationship] (contains) at ($(playlist)!0.5!(song)$) {contains};

    % Connections with cardinality
    \draw[thick] (artist) -- node[above] {1} (creates);
    \draw[thick] (creates) -- node[above] {N} (song);
    \draw[thick] (user) -- node[above] {1} (owns);
    \draw[thick] (owns) -- node[above] {N} (playlist);
    \draw[thick] (playlist) -- node[left] {M} (contains);
    \draw[thick] (contains) -- node[right] {N} (song);

\end{tikzpicture}
\end{center}

\begin{warningbox}[Συνήθες Λάθος]
Πολλοί φοιτητές ξεχνούν ότι η σχέση \texttt{contains} είναι M:N και χρειάζεται
ενδιάμεσο πίνακα στο σχεσιακό μοντέλο!
\end{warningbox}

\textbf{Τελική Απάντηση:} Το E-R διάγραμμα περιλαμβάνει 4 οντότητες και 3 συσχετίσεις
(1:N, 1:N, M:N).

% -----------------------------------------------------------------------------
% Exercise 1.2
% -----------------------------------------------------------------------------
\section{Άσκηση 1.2: Μετατροπή E-R σε Σχεσιακό (M:N)}

\begin{formulabox}[Εκφώνηση]
Μετατρέψτε την M:N σχέση \texttt{contains} μεταξύ \texttt{Playlist} και \texttt{Song}
σε σχεσιακούς πίνακες. Η σχέση έχει επιπλέον γνώρισμα \texttt{position} που δηλώνει
τη θέση του τραγουδιού στη playlist.
\end{formulabox}

\subsection*{Λύση}

\textbf{Βήμα 1: Κανόνας Μετατροπής M:N}

\textit{Γιατί:} Οι M:N σχέσεις δεν μπορούν να αναπαρασταθούν με ξένα κλειδιά σε έναν
από τους δύο πίνακες. Χρειάζεται ενδιάμεσος πίνακας.

\textbf{Βήμα 2: Δημιουργία Πινάκων}

\begin{lstlisting}[language=SQL]
-- Πίνακας Playlist
CREATE TABLE Playlist (
    playlist_id INT PRIMARY KEY,
    name VARCHAR(100),
    created_date DATE,
    user_id INT REFERENCES User(user_id)
);

-- Πίνακας Song
CREATE TABLE Song (
    song_id INT PRIMARY KEY,
    title VARCHAR(200),
    duration INT,
    release_year INT,
    artist_id INT REFERENCES Artist(artist_id)
);

-- Ενδιάμεσος πίνακας για M:N
CREATE TABLE Playlist_Song (
    playlist_id INT REFERENCES Playlist(playlist_id),
    song_id INT REFERENCES Song(song_id),
    position INT,  -- Γνώρισμα της σχέσης
    PRIMARY KEY (playlist_id, song_id)
);
\end{lstlisting}

\begin{infobox}[Σημαντικό]
Το πρωτεύον κλειδί του ενδιάμεσου πίνακα είναι το συνδυασμένο κλειδί
\texttt{(playlist\_id, song\_id)} που αποτρέπει διπλότυπες εγγραφές.
\end{infobox}

\textbf{Τελική Απάντηση:} Δημιουργούνται 3 πίνακες: \texttt{Playlist}, \texttt{Song},
και \texttt{Playlist\_Song} (ενδιάμεσος).

% -----------------------------------------------------------------------------
% Exercise 1.3
% -----------------------------------------------------------------------------
\section{Άσκηση 1.3: Weak Entity Μετατροπή}

\begin{formulabox}[Εκφώνηση]
Ένα σύστημα διαχείρισης παραγγελιών έχει:
\begin{itemize}
    \item \texttt{Order} με \underline{order\_id} και order\_date
    \item \texttt{OrderItem} (weak entity) με \underline{item\_number}, quantity, price
\end{itemize}
Το \texttt{OrderItem} εξαρτάται από το \texttt{Order} για την αναγνώρισή του.
Μετατρέψτε σε σχεσιακό μοντέλο.
\end{formulabox}

\subsection*{Λύση}

\textbf{Βήμα 1: Αναγνώριση Weak Entity}

\textit{Γιατί:} Μια weak entity δεν έχει αρκετά γνωρίσματα για να ταυτοποιηθεί μόνη
της. Το \texttt{item\_number} (π.χ. 1, 2, 3) επαναλαμβάνεται μεταξύ διαφορετικών
παραγγελιών.

\textbf{Βήμα 2: Σύνθετο Πρωτεύον Κλειδί}

Το πρωτεύον κλειδί της weak entity σχηματίζεται από:
\begin{itemize}
    \item Το πρωτεύον κλειδί της «ιδιοκτήτριας» οντότητας (owner entity)
    \item Το μερικό κλειδί (partial key) της weak entity
\end{itemize}

\textbf{Βήμα 3: Σχεσιακοί Πίνακες}

\begin{lstlisting}[language=SQL]
CREATE TABLE Order (
    order_id INT PRIMARY KEY,
    order_date DATE
);

CREATE TABLE OrderItem (
    order_id INT REFERENCES Order(order_id),
    item_number INT,
    quantity INT,
    price DECIMAL(10,2),
    PRIMARY KEY (order_id, item_number)  -- Σύνθετο κλειδί
);
\end{lstlisting}

\begin{warningbox}[Συνήθες Λάθος]
Μην ξεχνάτε: το ξένο κλειδί (\texttt{order\_id}) γίνεται \textbf{μέρος} του
πρωτεύοντος κλειδιού στη weak entity, όχι απλά ξένο κλειδί!
\end{warningbox}

\textbf{Τελική Απάντηση:} PK του OrderItem = \texttt{(order\_id, item\_number)}

% -----------------------------------------------------------------------------
% Exercise 1.4
% -----------------------------------------------------------------------------
\section{Άσκηση 1.4: Cardinality Constraints}

\begin{formulabox}[Εκφώνηση]
Για τις παρακάτω περιγραφές, προσδιορίστε το cardinality (1:1, 1:N, M:N):

\begin{enumerate}
    \item Ένας φοιτητής εγγράφεται σε πολλά μαθήματα, ένα μάθημα έχει πολλούς φοιτητές
    \item Κάθε υπάλληλος έχει ακριβώς έναν διευθυντή, ένας διευθυντής διευθύνει πολλούς υπαλλήλους
    \item Κάθε χώρα έχει μία πρωτεύουσα, κάθε πρωτεύουσα ανήκει σε μία χώρα
    \item Ένας συγγραφέας γράφει πολλά βιβλία, ένα βιβλίο μπορεί να έχει πολλούς συγγραφείς
\end{enumerate}
\end{formulabox}

\subsection*{Λύση}

\textbf{Μέθοδος:} Για κάθε περίπτωση ρωτάμε:
\begin{itemize}
    \item Πόσες οντότητες Β συνδέονται με μία οντότητα Α;
    \item Πόσες οντότητες Α συνδέονται με μία οντότητα Β;
\end{itemize}

\begin{center}
\begin{tabular}{clcl}
\toprule
\textbf{\#} & \textbf{Σχέση} & \textbf{Cardinality} & \textbf{Εξήγηση} \\
\midrule
1 & Student-Course & \textbf{M:N} & Πολλοί-σε-πολλούς \\
2 & Employee-Manager & \textbf{N:1} & Πολλοί υπάλληλοι, ένας διευθυντής \\
3 & Country-Capital & \textbf{1:1} & Ένα-προς-ένα \\
4 & Author-Book & \textbf{M:N} & Πολλοί-σε-πολλούς \\
\bottomrule
\end{tabular}
\end{center}

\begin{infobox}[Υλοποίηση σε Σχεσιακό]
\begin{itemize}
    \item \textbf{1:1}: FK σε οποιονδήποτε πίνακα (+ UNIQUE)
    \item \textbf{1:N}: FK στην πλευρά του N
    \item \textbf{M:N}: Ενδιάμεσος πίνακας (junction table)
\end{itemize}
\end{infobox}

% -----------------------------------------------------------------------------
% Exercise 1.5
% -----------------------------------------------------------------------------
\section{Άσκηση 1.5: Σύνθετο Schema (Hospital System)}

\begin{formulabox}[Εκφώνηση]
Σχεδιάστε E-R και μετατρέψτε σε σχεσιακό για νοσοκομειακό σύστημα:
\begin{itemize}
    \item \texttt{Doctor}: doctor\_id, name, specialty
    \item \texttt{Patient}: patient\_id, name, birth\_date
    \item \texttt{Appointment}: date, time, diagnosis (σχέση μεταξύ Doctor-Patient)
    \item Ένας γιατρός έχει πολλά ραντεβού, ένας ασθενής έχει πολλά ραντεβού
    \item Ένα ραντεβού συνδέει ακριβώς έναν γιατρό με έναν ασθενή
\end{itemize}
\end{formulabox}

\subsection*{Λύση}

\textbf{Βήμα 1: Ανάλυση}

Η σχέση Doctor-Patient είναι M:N (ένας γιατρός βλέπει πολλούς ασθενείς, ένας ασθενής
επισκέπτεται πολλούς γιατρούς), αλλά κάθε \textit{συγκεκριμένο} ραντεβού αφορά
ακριβώς έναν γιατρό και έναν ασθενή.

\textbf{Βήμα 2: E-R Διάγραμμα}

\begin{center}
\begin{tikzpicture}[node distance=4cm,
    entity/.style={rectangle, draw=primary, fill=box-blue, minimum width=2cm, minimum height=1cm, font=\bfseries},
    relationship/.style={diamond, draw=secondary, fill=box-yellow, aspect=2.5, font=\small}]

    \node[entity] (doctor) {Doctor};
    \node[entity, right=5cm of doctor] (patient) {Patient};
    \node[relationship] (appointment) at ($(doctor)!0.5!(patient)$) {Appointment};

    \draw[thick] (doctor) -- node[above] {1} (appointment);
    \draw[thick] (appointment) -- node[above] {N} (patient);
    \draw[thick] (patient) -- node[below] {1} (appointment);
    \draw[thick] (appointment) -- node[below] {M} (doctor);

    % Attributes of relationship
    \node[ellipse, draw=gray, fill=light-gray, font=\small, above=0.5cm of appointment] (date) {date, time};
    \node[ellipse, draw=gray, fill=light-gray, font=\small, below=0.5cm of appointment] (diag) {diagnosis};
    \draw[dashed] (appointment) -- (date);
    \draw[dashed] (appointment) -- (diag);
\end{tikzpicture}
\end{center}

\textbf{Βήμα 3: Σχεσιακό Μοντέλο}

\begin{lstlisting}[language=SQL]
CREATE TABLE Doctor (
    doctor_id INT PRIMARY KEY,
    name VARCHAR(100),
    specialty VARCHAR(50)
);

CREATE TABLE Patient (
    patient_id INT PRIMARY KEY,
    name VARCHAR(100),
    birth_date DATE
);

CREATE TABLE Appointment (
    appointment_id INT PRIMARY KEY,
    doctor_id INT REFERENCES Doctor(doctor_id),
    patient_id INT REFERENCES Patient(patient_id),
    appointment_date DATE,
    appointment_time TIME,
    diagnosis TEXT
);
\end{lstlisting}

\textbf{Τελική Απάντηση:} 3 πίνακες με τον \texttt{Appointment} να περιέχει FKs και
προς τους δύο.
